% ---------------------------------------------------------------------------
% The change of variables (x1, x2) -> (ln tau, Y): the allowed triangle, the
% lines of constant tau and Y, and the little rectangle whose image has area
% tau d ln tau dY. Numbers (tau_min, our event) from section [1] and [2] of the
% -v log of standalone_qqbar_bbbar_xsec.py.
% ---------------------------------------------------------------------------
\begin{tikzpicture}[x=1cm,y=1cm,
  every node/.style={font=\footnotesize},
  axis/.style={-{Stealth},thick},
  iso/.style={gray!60,thin},
  lab/.style={font=\scriptsize,align=center},
  head/.style={font=\small\bfseries},
  ev/.style={star,star points=5,star point ratio=2.2,fill=red!80!black,draw=white,inner sep=0pt,minimum size=3.2mm},
  bigarrow/.style={-{Stealth[length=3mm,width=3mm]},gray!60,line width=2.4pt}]
% ---------------- (a) the (ln x1, ln x2) plane: a triangle ---------------------------------
% horizontal: -ln x1 (x1 decreases to the right), vertical: -ln x2; origin = the tau = 1 corner
\begin{scope}
\node[head,anchor=west] at (-0.2,4.9) {(a) the $(\ln x_1,\ln x_2)$ plane};
\fill[blue!8] (0,0) -- (4.2,0) -- (0,4.2) -- cycle;
\draw[axis] (-0.3,0) -- (4.9,0) node[below left] {$-\ln x_1$};
\draw[axis] (0,-0.3) -- (0,4.9) node[below left] {$-\ln x_2$};
\node[lab,anchor=north] at (4.2,-0.1) {$x_1=\tau_{\min}$};
\node[lab,anchor=east] at (-0.1,4.2) {$x_2=\tau_{\min}$};
\node[lab,anchor=north east] at (-0.05,-0.05) {$x_1=x_2=1$\\$(\tau=1)$};
\draw[thick] (4.2,0) -- (0,4.2);
\node[lab,rotate=-45] at (2.6,2.0) {$x_1x_2=\tau_{\min}=2.17\times10^{-4}$};
\foreach \t in {1.4,2.8} {\draw[iso] (\t,0) -- (0,\t);}
\node[lab,rotate=-45,text=gray!80!black] at (1.75,1.75) {$\tau=\mathrm{const}$};
\draw[iso,dashed] (0,0) -- (2.1,2.1);
\draw[iso,dashed] (1.4,0) -- (2.8,1.4);
\draw[iso,dashed] (0,1.4) -- (1.4,2.8);
\node[lab,rotate=45,text=gray!80!black] at (1.3,1.55) {$Y=0$};
\node[lab,rotate=45,text=gray!80!black] at (2.15,0.95) {$Y<0$};
\node[lab,rotate=45,text=gray!80!black] at (0.75,2.35) {$Y>0$};
\node[ev] at (3.05,0.47) {};
\node[lab,anchor=east,text=red!80!black] at (2.9,0.5) {our event, $Y=-2.59$};
\node[lab,text=black!65,text width=4.6cm,anchor=north west] at (-0.2,-0.95) {The region above the threshold line is a triangle: $\tau$ is constant along the anti-diagonals, $Y$ along the diagonals ($Y<0$: $x_1<x_2$, our side).};
\end{scope}
\draw[bigarrow] (5.3,2.2) -- (6.1,2.2);
% ---------------- (b) the (ln tau, Y) plane ---------------------------------------------------
\begin{scope}[xshift=7.2cm]
\node[head,anchor=west] at (-0.2,4.9) {(b) the $(\ln\tau,Y)$ plane};
\fill[blue!8] (0,2.2) -- (4.2,4.3) -- (4.2,0.1) -- cycle;
\draw[axis] (-0.3,2.2) -- (4.9,2.2) node[below left] {$\ln\tau$};
\draw[axis] (0,-0.3) -- (0,4.9) node[below left] {$Y$};
\node[lab,anchor=north] at (0,2.1) {$\ln\tau_{\min}$};
\node[lab,anchor=north] at (4.2,2.1) {$0$};
\draw[thick] (0,2.2) -- (4.2,4.3) node[lab,pos=0.55,above,sloped] {$Y_{\max}=-\tfrac12\ln\tau$};
\draw[thick] (0,2.2) -- (4.2,0.1) node[lab,pos=0.88,below,sloped] {$-Y_{\max}$};
\foreach \t in {1.4,2.8} {\draw[iso] (\t,{2.2-(4.2-\t)/2}) -- (\t,{2.2+(4.2-\t)/2});}
\foreach \y in {1.5,2.9} {\draw[iso,dashed] ({4.2-2*abs(\y-2.2)},\y) -- (4.2,\y);}
\draw[iso,dashed] (0,2.2) -- (4.2,2.2);
% the little rectangle
\fill[orange!60] (1.55,1.5) rectangle (1.95,1.9);
\draw[thick,orange!80!black] (1.55,1.5) rectangle (1.95,1.9);
\node[lab,anchor=south west,text=orange!80!black] at (1.9,1.9) {$\mathrm{d}\ln\tau\,\mathrm{d}Y$};
\node[ev] at (0.68,0.95) {};
\draw[red!80!black,thin] (0.68,0.95) -- (1.6,-0.15);
\node[lab,anchor=north,text=red!80!black] at (1.9,-0.15) {our event: $\ln\tau=-7.07$, $Y=-2.59$ ($73\%$ of $Y_{\max}=3.54$)};
\node[lab,text=black!65,text width=4.6cm,anchor=north west] at (-0.2,-1.15) {The unit cube is mapped linearly onto this triangle: $u_1$ along $\ln\tau$ over the length $\ln(1/\tau_{\min})=8.44$, $u_2$ along $Y$ over $2Y_{\max}(\tau)$.};
\end{scope}
\draw[bigarrow] (12.5,2.2) -- (13.3,2.2);
% ---------------- (c) the Jacobian in the linear (x1, x2) plane --------------------------------
\begin{scope}[xshift=14.4cm]
\node[head,anchor=west] at (-0.2,4.9) {(c) the Jacobian, $\mathrm{d}x_1\mathrm{d}x_2=\tau\,\mathrm{d}\ln\tau\,\mathrm{d}Y$};
\draw[axis] (-0.3,0) -- (4.9,0) node[below left] {$x_1$};
\draw[axis] (0,-0.3) -- (0,4.9) node[below left] {$x_2$};
% a hyperbola tau = const and its neighbour
\draw[iso] plot[domain=0.5:4.3,samples=60] (\x,{2.0/\x});
\draw[iso] plot[domain=0.62:4.3,samples=60] (\x,{2.6/\x});
\node[lab,text=gray!80!black,anchor=west] at (3.4,0.95) {$\tau$ and $\tau e^{\mathrm{d}\ln\tau}$};
% rays Y = const and neighbour
\draw[iso,dashed] (0,0) -- (3.9,3.9);
\draw[iso,dashed] (0,0) -- (4.3,2.9);
\node[lab,text=gray!80!black,anchor=west,rotate=34] at (3.2,2.15) {$Y$ and $Y+\mathrm{d}Y$};
% the parallelogram: corners at the intersections
\coordinate (p1) at (1.414,1.414); \coordinate (p2) at (1.612,1.612);
\coordinate (p3) at (1.964,1.324); \coordinate (p4) at (1.723,1.161);
\fill[orange!60] (p1) -- (p2) -- (p3) -- (p4) -- cycle;
\draw[thick,orange!80!black] (p1) -- (p2) -- (p3) -- (p4) -- cycle;
\node[lab,anchor=south west,text=orange!80!black] at (2.0,1.55) {area $=\tau\,\mathrm{d}\ln\tau\,\mathrm{d}Y$};
\node[lab,text=black!65,text width=4.6cm,anchor=north west] at (-0.2,-0.95) {The rectangle becomes a tilted parallelogram $\tau$ times larger. Since the libraries return $xf$, the two factors of $\tau$ cancel: at our event $f_{\bar u}f_u\,\tau=0.2413=[x_1f_{\bar u}][x_2f_u]$.};
\end{scope}
\end{tikzpicture}
